The second-moment bound gives tightness. Higher even moments give uniform integrability when passing moments to subsequential limits. The Gaussian moment sequence is determinate, for example by Carleman\textquotesingle s criterion, since \(((2r-1)!!)^{1/(2r)}\leq\sqrt{2r}\). Every subsequential limit is standard normal. Convergence of the second moments upgrades weak convergence to quadratic Wasserstein convergence. End of proof.

\subsection{Critical and subcritical propagator spectra}\label{floquet:critical-and-subcritical-propagator-spectra}

For \(\varepsilon=\tau/\sqrt{2K}\), (FN.17) becomes \[
 \left|\tau_-(U_\varepsilon^{mK})-
              \tau_-(e^{im\tau X_K})\right|
       \leq \frac{6|m|\tau^2}{\sqrt K}.                        \tag{FN.21}
\] Lemma 2 and bounded continuity of exp(im \(\tau\) x) give (FN.A). Integer Fourier coefficients determine measures on the circle: trigonometric polynomials are dense in continuous functions there. Applying the characteristic function of a normal variable proves (FN.B), and taking absolute squares proves (FN.C). No eigenphase labelling across the old degenerate sectors is used or asserted.

There is also a genuine extension of the earlier subcritical global transfer, but for \textbf{folded phases} rather than individually isolated clusters. Let \(\varepsilon_{K}\) be nonzero and \(\tau_K=\varepsilon_K\sqrt{2K}\) tend to zero. If \(e^{i\theta_l}\) are the eigenvalues of \(U_{\varepsilon,K}\), define \[
 y_l=\frac{\operatorname{Arg}(e^{iK\theta_l})}{\tau_K},
 \qquad \operatorname{Arg}\in[-\pi,\pi).
\]

\textbf{Corollary B (subcritical folded Gaussian law).} The empirical measures of \(y_{l}\) converge in quadratic Wasserstein distance to the standard normal law whenever \(\varepsilon_K\sqrt K\) tends to zero.

\textbf{Proof with a quantitative tail control.} For two unitary matrices, the squared absolute overlaps of their orthonormal eigenbases form a doubly stochastic matrix. The corresponding average squared chord distance equals their normalized squared Hilbert--Schmidt difference. Decomposing this matrix into permutation matrices gives an eigenvalue matching with no larger average cost. The decomposition can be obtained by Hall's matching criterion and successive subtraction; it uses no eigenvalue simplicity. Thus (FN.16), telescoped K times, matches the powered spectrum to \(e^{i\tau_Kx_l}\), where \(x_{l}\) are the eigenvalues of \(X_{K}\), with root mean square chord error at most \(6\tau_K^2/\sqrt K\).

If \(|\tau_Kx_l|\leq\pi/2\), principal argument differences obey \(|\operatorname{Arg}u-\tau_Kx_l|\leq4|u-e^{i\tau_Kx_l}|\). On the complementary set, \(|\operatorname{Arg}u/\tau_K-x_l|\leq3|x_l|\). Consequently \[
 \begin{split}
 W_2(\operatorname{law}(y),\operatorname{law}(X_K))
 &\leq\frac{24|\tau_K|}{\sqrt K}
   +3\sqrt{\tau_-\!\left(X_K^2
       \mathbf1_{|X_K|>\pi/(2|\tau_K|)}\right)}\\
 &\leq\frac{24|\tau_K|}{\sqrt K}
       +\frac{6|\tau_K|}{\pi}\sqrt{\tau_-(X_K^4)}.
 \end{split}\tag{FN.22}
\] The fourth moment stays bounded by Lemma 2, so the right side tends to zero. Combine with the \(W_{2}\) conclusion of that lemma. End of proof.

The former isolated-cluster argument required \(\varepsilon_KK^{3/2}\) tending to zero. Corollary B permits an additional factor K in the kick scale, but its observable is a folded spectral measure. It must not be restated as an individual-cluster rank or no-overlap theorem.

\subsection{A compact-test microscopic obstruction}\label{floquet:a-compact-test-microscopic-obstruction}

A negative statement about microscopic correlations needs less information than a positive universality theorem. The recurrence law above supplies such a statement through Fourier positivity. This section concerns the whole odd spectrum, with its multiplicities and the \textbf{global mean-density unfolding}. It is not a claim about a further symmetry sector or an unfolding which first removes a nonconstant mesoscopic intensity profile.

Let \(x_{1}\),\ldots,\(x_{d}\) be the eigenphases of a d-dimensional unitary U, measured in turns, so its eigenvalues are \(e^{2\pi ix_j}\). For a compactly supported continuous test \(\phi\) define the ordered factorial pair statistic \[
 C_d(U;\phi)=\frac1d\sum_{a\ne b}\sum_{n\in\mathbb Z}
              \phi\bigl(d(x_a-x_b+n)\bigr).                    \tag{FN.23}
\] Distinct indices, not distinct eigenvalues, are required in this definition. For \(L>0\) use the compact triangular test \[
 \phi_L(s)=\frac1L(1-|s|/L)_+,\quad
 \int\phi_L=1,\quad \phi_L(0)=1/L,\quad
 \widehat\phi_L(u)=\operatorname{sinc}^2(Lu)\geq0,               \tag{FN.24}
\] where \(\operatorname{sinc}(v)=\sin(\pi v)/(\pi v)\), with value 1 at zero. Fourier transforms use \(e^{-2\pi ius}\). Absolute convergence of the periodized Fourier series (the coefficients decay quadratically) gives the exact identity \[
 C_d(U;\phi_L)=d^{-2}\sum_{t\in\mathbb Z}
   \widehat\phi_L(t/d)(|\operatorname{Tr}U^t|^2-d).              \tag{FN.25}
\] If \(d>L\), the total diagonal subtraction is exactly 1/L, because \(\sum_n\phi_L(dn)=\phi_L(0)\). For any distinct positive integers \(r_{1}\),\ldots,\(r_{M}\), positivity therefore gives \[
 C_d(U;\phi_L)\geq1-\frac1L+
   2\sum_{j=1}^M\operatorname{sinc}^2(Lr_j/d)
                   |d^{-1}\operatorname{Tr}U^{r_j}|^2.         \tag{FN.26}
\] No control or truncation of the remaining high-frequency traces is needed for this lower bound; their contribution is nonnegative.

There is an equivalent counting interpretation. Let \(N_{L}\)(y) count the phases in the circular interval {[}y,y+L/d), with y uniform in {[}0,1). Then \(E_{y}\) \(N_{L}\)=L and an overlap-of-intervals calculation gives \[
 \operatorname{Var}_y N_L
      =L^2\bigl(C_d(U;\phi_L)-1\bigr)+L.                       \tag{FN.27}
\] This identity is also verified independently by integrating the piecewise-constant counting function, rather than reusing the pair sum.

\textbf{Theorem C (obstruction and necessary kick scale).} For the kicked Pell family of Theorem A, put \(\tau_K=\varepsilon_K\sqrt{2K}\). If \(\tau_{K}\) tends to a finite nonzero \(\tau\), then for every fixed \(L>0\), \[
 \liminf_{K\to\infty} C_{d_N}(U_{\varepsilon_K,N};\phi_L)
      \geq \Theta(\tau)-\frac1L,\qquad
 \Theta(\tau)=\sum_{m\in\mathbb Z}e^{-m^2\tau^2}>1,             \tag{FN.28}
\] and \[
 \liminf_{K\to\infty}\operatorname{Var}_y N_L
      \geq L^2\bigl(\Theta(\tau)-1\bigr).                      \tag{FN.29}
\] If \(\tau_K\to0\), then \(C_{d_N}(U_{\varepsilon_K,N};\phi_L)\to+\infty\) for every fixed \(L\). Consequently, convergence of the globally unfolded pair measures to a bounded density \(\rho(s)\) with \(\rho(s)\to1\) as \(|s|\to\infty\) requires \[
             |\varepsilon_K|\sqrt K\longrightarrow\infty.     \tag{FN.30}
\] This is a necessary condition only, not a sufficient universality criterion.

\textbf{Proof.} The recurrence trace limit remains valid when \(\tau_K\to\tau\) varies with \(K\). Lemma 2 bounds \(\tau_-(|X_K|)\), and \(|\tau_-(e^{im\tau_KX_K})-\tau_-(e^{im\tau X_K})|\leq |m|\,|\tau_K-\tau|\,\tau_-(|X_K|)\). The error in (FN.21) is uniform for bounded \(\tau_K\). Thus the normalized \(mK\)-trace tends to \(e^{-m^2\tau^2/2}\) for each fixed \(m\), including \(\tau=0\).

Apply (FN.26) with \(r_j=jK\). Since \(K/d_N\to0\), the sinc factors tend to one. At each fixed \(M\), the resulting lower bound is \(1-1/L+2\sum_{m=1}^M e^{-m^2\tau^2}\). Let \(M\) tend to infinity \emph{after} taking the lower limit in \(K\). This proves (FN.28), or divergence when \(\tau=0\). Equation (FN.27) gives (FN.29).

Suppose the pair measures converge vaguely to \(\rho\)(s)ds, where \(\rho\) is bounded and clusters to unit intensity. For every fixed L the limiting triangular statistic would be \[
 A_L=\int_{-1}^1(1-|u|)\rho(Lu)\,du.
\] Dominated convergence gives \(A_L\to1\) as \(L\to\infty\). But (FN.28) requires \(A_L\geq\Theta(\tau)-1/L\) for every \(L\), a contradiction. For \(\tau=0\), a finite locally bounded target already contradicts the fixed-\(L\) divergence. If \(|\varepsilon_K|\sqrt K\) does not tend to infinity, a bounded subsequence of \(\tau_K\) has a further subsequence converging to a finite limit. The same contradiction applies. End of proof.

